\pdfbookmark[1]{Hassler. 1831 März 31.}{hassler1831}\par\Needspace{10\baselineskip}\begin{center}\rule{0.7\linewidth}{0.4pt}\par Göttingische gelehrte Anzeigen. 1831 März 31.\par\rule{0.7\linewidth}{0.4pt}\end{center}

\noindent\textit{Tabulae logarithmicae et trigonometricae notis septem decimalibus expressae.}
\textit{In forma minima. Purgatae ab erroribus praecedentium tabularum cura} F. K. \textsc{Hassler.} New York. G. C. und H. Carvill.

Diese transatlantische Ausgabe der Logarithmentafeln zeichnet sich durch die ganz vorzügliche Sauberkeit ihres Stereotypdrucks aus. Sie enthält durchgehend auf sieben Dezimalstellen die Logarithmen der Zahlen bis 100000 sowie die Logarithmen der Sinus, Tangenten, Kosinus und Kotangenten im ersten Grad für jede Sekunde, in den beiden folgenden Graden in Abständen von zehn Sekunden und in allen übrigen Graden des Quadranten in Abständen von dreißig Sekunden. Dazu kommen die natürlichen Sinus und Tangenten durchgehend in Abständen von dreißig Sekunden. Und dies alles auf 312 Seiten in Kleinoktav oder Großduodez. Eine solche Verdichtung war freilich nur mit sehr kleinen Drucktypen zu erreichen. Bei aller Schönheit dürften sie für die meisten nicht kurzsichtigen Augen zum täglichen Gebrauch beinahe zu klein sein. Auf die Richtigkeit scheint ganz besondere Sorgfalt verwendet worden zu sein; wenigstens ist uns bei dem eine Zeit lang versuchten häufigen Gebrauch kein einziger Druckfehler aufgefallen.

